\documentclass[12pt]{article}
\usepackage[utf8]{inputenc}
\usepackage[T1]{fontenc}
\usepackage{hyperref}
\title{Mewati Language Resources for K-12 Curriculum Development}
\author{Compiled by Claude Sonnet 5.5 (claude-sonnet-5-5) for Andr\'as Kornai}
\date{}
\begin{document}
\maketitle

\section{Introduction}
Mewati (ISO 639-3 wtm; Glottocode mewa1250) is an Indo-European (Indo-Aryan) language of the Mewat region of India, and also of Pakistan. It is often treated as a Rajasthani dialect, and search results report that it has no standard script. Its literature has largely been passed on orally. This compilation rests on one round of web searches. No textbooks or dictionaries turned up.

\section{Primary Educational Resources}
I found no primer or textbook. The UNESCO language atlas entry mentions Mewati folksongs, stories and literature written in Devanagari on a Mewati website and in an Internet Archive scan: \url{https://ar.wal.unesco.org/index.php/en/countries/india/languages/mewati}. I did not see the underlying URLs, so those items are not listed in the json.

\section{Linguistic and Reference Materials}
\begin{itemize}
\item Overview: \url{https://wikipedia.com/wiki/ISO_639:wtm}
\item ``Morphology of Mewati Language'' (conference presentation): \url{https://www.cle.org.pk/clt16/Presentations/Morphology of Mewati Language.pdf}
\item Nadia Fareed, M.Phil thesis on the morphemic structure of Lahori Mewati (UMT Lahore, 2016): \url{https://escholar.umt.edu.pk/items/3cdfd632-7021-486b-9417-575d0ca1a4ae/full}
\item LINGUIST List announcement, issue 18.2197: \url{https://linguistlist.org/issues/18/2197/}
\end{itemize}
The search summary also named Lakhan Gusain's \emph{Mewati} (Lincom, 2003/2007) as the main grammar, and older dissertations by A.\ R.\ Fatihi (1987) and D.\ R.\ Bharadwaj (1974). I saw no URL for these, so they are not in the json. I found no Mewati dictionary.

\section{Religious and Community Resources}
A news story discusses how Mewati gained status as a separate language in Pakistan: \url{https://loksujag.com/story/How-Did-Mewati-Gain-Language-Status-76-Years-After-Migration-eng}. The search summary also mentioned a ``Mewati Dunya'' website, but I saw no URL for it.

\section{Digital and Audio Resources}
The UNESCO entry says folk tales, songs and video exist on the Internet Archive. I saw no specific item URLs.

\section{Curriculum Development Context}
Mewati is not known to be taught in schools, and it lacks a standard script, so a curriculum would first need a script decision. Devanagari is used in some existing community publications, and Pakistani Mewati writers have reportedly produced grammar books.

\end{document}